%%%%%%%%%%%%%%%%%%%%%%%%%%%%%%%%%%%%%%%%%%%%%%%%%%%%%%%%%%%%
% 11.6 MONTE CARLO METHODS
% The Composition of Advanced Mathematics
%%%%%%%%%%%%%%%%%%%%%%%%%%%%%%%%%%%%%%%%%%%%%%%%%%%%%%%%%%%%

\section{Monte Carlo Methods}
\label{sec:monte-carlo-methods}

\prerequisite{
Probability, random variables, expectation, variance, probability
distributions, laws of large numbers, central limit theorems, numerical
integration, statistics, stochastic processes, and basic programming.
}

\learningobjectives{
By the end of this section, the reader should be able to:
\begin{itemize}
    \item explain the basic Monte Carlo principle;
    \item estimate expectations using random samples;
    \item understand why Monte Carlo error decreases at rate \(N^{-1/2}\);
    \item distinguish bias from sampling variance;
    \item construct Monte Carlo estimators;
    \item estimate the standard error of a Monte Carlo estimate;
    \item construct approximate Monte Carlo confidence intervals;
    \item approximate definite integrals using random sampling;
    \item approximate multidimensional integrals;
    \item recognize the dimensional advantages of Monte Carlo integration;
    \item simulate common probability distributions;
    \item understand pseudorandom number generation;
    \item distinguish pseudorandom and truly random sequences;
    \item understand reproducibility through random seeds;
    \item apply inverse-transform sampling;
    \item apply acceptance--rejection sampling;
    \item recognize transformation methods for simulation;
    \item understand importance sampling;
    \item derive importance-sampling estimators;
    \item understand variance reduction;
    \item apply antithetic variates conceptually;
    \item apply control variates;
    \item understand stratified sampling;
    \item recognize Latin hypercube sampling conceptually;
    \item understand common random numbers conceptually;
    \item estimate rare-event probabilities;
    \item understand why naive Monte Carlo is inefficient for rare events;
    \item recognize sequential Monte Carlo conceptually;
    \item understand Markov chain Monte Carlo;
    \item identify stationary and target distributions;
    \item understand the Metropolis--Hastings algorithm;
    \item understand Gibbs sampling conceptually;
    \item distinguish independent Monte Carlo from MCMC;
    \item understand burn-in and initialization issues;
    \item recognize autocorrelation in Markov-chain samples;
    \item understand effective sample size conceptually;
    \item understand Monte Carlo diagnostics;
    \item recognize quasi-Monte Carlo methods;
    \item understand low-discrepancy sequences conceptually;
    \item distinguish Monte Carlo from quasi-Monte Carlo;
    \item understand bootstrap simulation conceptually;
    \item recognize Monte Carlo methods for stochastic differential
          equations;
    \item estimate expectations from simulated SDE paths;
    \item understand nested simulation;
    \item recognize uncertainty quantification;
    \item perform convergence studies for stochastic simulation;
    \item connect Monte Carlo methods with probability, statistics,
          optimization, finance, physics, engineering, Bayesian inference,
          and scientific computing.
\end{itemize}
}

%%%%%%%%%%%%%%%%%%%%%%%%%%%%%%%%%%%%%%%%%%%%%%%%%%%%%%%%%%%%
\subsection{What Is a Monte Carlo Method?}
\label{subsec:what-is-monte-carlo}
%%%%%%%%%%%%%%%%%%%%%%%%%%%%%%%%%%%%%%%%%%%%%%%%%%%%%%%%%%%%

Monte Carlo methods use random sampling to approximate mathematical
quantities.

The basic idea is:

\[
\boxed{
\text{random samples}
\longrightarrow
\text{sample average}
\longrightarrow
\text{approximation}.
}
\]

Monte Carlo methods are especially useful when direct deterministic
computation is difficult.

%%%%%%%%%%%%%%%%%%%%%%%%%%%%%%%%%%%%%%%%%%%%%%%%%%%%%%%%%%%%
\subsection{Expectation as the Central Quantity}
\label{subsec:monte-carlo-expectation}
%%%%%%%%%%%%%%%%%%%%%%%%%%%%%%%%%%%%%%%%%%%%%%%%%%%%%%%%%%%%

Suppose

\[
X
\]

is a random variable and we want to compute

\[
\boxed{
\mu
=
\mathbb{E}[g(X)].
}
\]

If

\[
X_1,\ldots,X_N
\]

are independent samples from the distribution of \(X\), define

\[
\boxed{
\widehat{\mu}_N
=
\frac1N
\sum_{i=1}^{N}
g(X_i).
}
\]

This is the standard Monte Carlo estimator.

%%%%%%%%%%%%%%%%%%%%%%%%%%%%%%%%%%%%%%%%%%%%%%%%%%%%%%%%%%%%
\subsection{Unbiasedness}
\label{subsec:monte-carlo-unbiasedness}
%%%%%%%%%%%%%%%%%%%%%%%%%%%%%%%%%%%%%%%%%%%%%%%%%%%%%%%%%%%%

If

\[
\mathbb{E}[|g(X)|]<\infty,
\]

then

\[
\mathbb{E}
[
\widehat{\mu}_N
]
=
\frac1N
\sum_{i=1}^{N}
\mathbb{E}[g(X_i)].
\]

Since the samples have identical distributions,

\[
\boxed{
\mathbb{E}
[
\widehat{\mu}_N
]
=
\mu.
}
\]

Thus the ordinary sample-average Monte Carlo estimator is unbiased.

%%%%%%%%%%%%%%%%%%%%%%%%%%%%%%%%%%%%%%%%%%%%%%%%%%%%%%%%%%%%
\subsection{Variance of the Monte Carlo Estimator}
\label{subsec:mc-estimator-variance}
%%%%%%%%%%%%%%%%%%%%%%%%%%%%%%%%%%%%%%%%%%%%%%%%%%%%%%%%%%%%

Assume

\[
\operatorname{Var}(g(X))
=
\sigma^2
<
\infty.
\]

For independent samples,

\[
\operatorname{Var}
(
\widehat{\mu}_N
)
=
\frac1{N^2}
\sum_{i=1}^{N}
\sigma^2.
\]

Therefore,

\[
\boxed{
\operatorname{Var}
(
\widehat{\mu}_N
)
=
\frac{
\sigma^2
}{
N
}.
}
\]

The standard deviation is

\[
\boxed{
\frac{\sigma}{\sqrt{N}}.
}
\]

%%%%%%%%%%%%%%%%%%%%%%%%%%%%%%%%%%%%%%%%%%%%%%%%%%%%%%%%%%%%
\subsection{Monte Carlo Convergence Rate}
\label{subsec:mc-convergence-rate}
%%%%%%%%%%%%%%%%%%%%%%%%%%%%%%%%%%%%%%%%%%%%%%%%%%%%%%%%%%%%

The characteristic Monte Carlo error scale is

\[
\boxed{
O(N^{-1/2}).
}
\]

To reduce standard error by a factor of \(2\), the sample size must be
multiplied by approximately

\[
\boxed{
4.
}
\]

To reduce error by a factor of \(10\), one generally needs about

\[
\boxed{
100
}
\]

times as many independent samples.

%%%%%%%%%%%%%%%%%%%%%%%%%%%%%%%%%%%%%%%%%%%%%%%%%%%%%%%%%%%%
\subsection{Law of Large Numbers}
\label{subsec:mc-law-large-numbers}
%%%%%%%%%%%%%%%%%%%%%%%%%%%%%%%%%%%%%%%%%%%%%%%%%%%%%%%%%%%%

Under standard assumptions, the Law of Large Numbers implies

\[
\boxed{
\widehat{\mu}_N
\to
\mu
}
\]

as

\[
N\to\infty.
\]

Thus the sample average converges to the desired expectation.

The Law of Large Numbers gives the basic consistency of Monte Carlo
estimation.

%%%%%%%%%%%%%%%%%%%%%%%%%%%%%%%%%%%%%%%%%%%%%%%%%%%%%%%%%%%%
\subsection{Central Limit Theorem}
\label{subsec:mc-central-limit-theorem}
%%%%%%%%%%%%%%%%%%%%%%%%%%%%%%%%%%%%%%%%%%%%%%%%%%%%%%%%%%%%

Under suitable finite-variance assumptions,

\[
\boxed{
\sqrt{N}
\left(
\widehat{\mu}_N-\mu
\right)
\overset{d}{\longrightarrow}
N(0,\sigma^2).
}
\]

Therefore, for large \(N\),

\[
\widehat{\mu}_N
\]

is approximately normally distributed around \(\mu\).

This allows approximate confidence intervals.

%%%%%%%%%%%%%%%%%%%%%%%%%%%%%%%%%%%%%%%%%%%%%%%%%%%%%%%%%%%%
\subsection{Monte Carlo Standard Error}
\label{subsec:mc-standard-error}
%%%%%%%%%%%%%%%%%%%%%%%%%%%%%%%%%%%%%%%%%%%%%%%%%%%%%%%%%%%%

The population variance \(\sigma^2\) is usually unknown.

Estimate it using

\[
\boxed{
s^2
=
\frac1{N-1}
\sum_{i=1}^{N}
\left(
g(X_i)-\widehat{\mu}_N
\right)^2.
}
\]

Then the estimated Monte Carlo standard error is

\[
\boxed{
\widehat{\operatorname{SE}}
(
\widehat{\mu}_N
)
=
\frac{s}{\sqrt{N}}.
}
\]

%%%%%%%%%%%%%%%%%%%%%%%%%%%%%%%%%%%%%%%%%%%%%%%%%%%%%%%%%%%%
\subsection{Approximate Confidence Interval}
\label{subsec:mc-confidence-interval}
%%%%%%%%%%%%%%%%%%%%%%%%%%%%%%%%%%%%%%%%%%%%%%%%%%%%%%%%%%%%

For sufficiently large independent samples, an approximate confidence
interval is

\[
\boxed{
\widehat{\mu}_N
\pm
z_{1-\alpha/2}
\frac{s}{\sqrt{N}}.
}
\]

For an approximate \(95\%\) interval,

\[
z_{0.975}
\approx1.96.
\]

Thus,

\[
\boxed{
\widehat{\mu}_N
\pm
1.96
\frac{s}{\sqrt{N}}.
}
\]

%%%%%%%%%%%%%%%%%%%%%%%%%%%%%%%%%%%%%%%%%%%%%%%%%%%%%%%%%%%%
\subsection{Worked Example: Monte Carlo Mean}
\label{subsec:worked-mc-mean}
%%%%%%%%%%%%%%%%%%%%%%%%%%%%%%%%%%%%%%%%%%%%%%%%%%%%%%%%%%%%

\begin{workedexamplebox}{Estimating an Expectation}

Suppose the simulated values of \(g(X)\) are

\[
2,\quad4,\quad3,\quad5,\quad6.
\]

Then

\[
\widehat{\mu}_5
=
\frac{
2+4+3+5+6
}{
5
}.
\]

Therefore,

\[
\widehat{\mu}_5
=
4.
\]

\begin{answerbox}
\[
\boxed{
\widehat{\mu}_5=4.
}
\]
\end{answerbox}

The estimate would become more reliable as the number of independent
samples increases.

\end{workedexamplebox}

%%%%%%%%%%%%%%%%%%%%%%%%%%%%%%%%%%%%%%%%%%%%%%%%%%%%%%%%%%%%
\subsection{Bias and Variance}
\label{subsec:mc-bias-variance}
%%%%%%%%%%%%%%%%%%%%%%%%%%%%%%%%%%%%%%%%%%%%%%%%%%%%%%%%%%%%

The mean-squared error satisfies

\[
\boxed{
\operatorname{MSE}
(
\widehat{\theta}
)
=
\operatorname{Var}
(
\widehat{\theta}
)
+
\operatorname{Bias}
(
\widehat{\theta}
)^2.
}
\]

Standard sample-average Monte Carlo is often unbiased.

Other simulation methods may introduce bias in exchange for lower variance
or computational efficiency.

Examples include:

\[
\boxed{
\text{time discretization of SDEs},
}
\]

and

\[
\boxed{
\text{finite burn-in or approximation in MCMC}.
}
\]

%%%%%%%%%%%%%%%%%%%%%%%%%%%%%%%%%%%%%%%%%%%%%%%%%%%%%%%%%%%%
\subsection{Monte Carlo Integration}
\label{subsec:monte-carlo-integration}
%%%%%%%%%%%%%%%%%%%%%%%%%%%%%%%%%%%%%%%%%%%%%%%%%%%%%%%%%%%%

Suppose

\[
I
=
\int_a^b
f(x)\,dx.
\]

Let

\[
U
\sim
\operatorname{Uniform}(a,b).
\]

Then

\[
\mathbb{E}[f(U)]
=
\frac1{b-a}
\int_a^b
f(x)\,dx.
\]

Therefore,

\[
\boxed{
I
=
(b-a)
\mathbb{E}[f(U)].
}
\]

A Monte Carlo estimator is

\[
\boxed{
\widehat{I}_N
=
\frac{
b-a
}{
N
}
\sum_{i=1}^{N}
f(U_i).
}
\]

%%%%%%%%%%%%%%%%%%%%%%%%%%%%%%%%%%%%%%%%%%%%%%%%%%%%%%%%%%%%
\subsection{Worked Example: Monte Carlo Integral}
\label{subsec:worked-mc-integral}
%%%%%%%%%%%%%%%%%%%%%%%%%%%%%%%%%%%%%%%%%%%%%%%%%%%%%%%%%%%%

\begin{workedexamplebox}{Estimating \(\int_0^1 x^2\,dx\)}

Let

\[
U\sim\operatorname{Uniform}(0,1).
\]

Then

\[
\int_0^1x^2\,dx
=
\mathbb{E}[U^2].
\]

Therefore,

\[
\widehat{I}_N
=
\frac1N
\sum_{i=1}^{N}
U_i^2.
\]

The exact value is

\[
\frac13.
\]

\begin{answerbox}
\[
\boxed{
\widehat{I}_N
=
\frac1N
\sum_{i=1}^{N}
U_i^2
\to
\frac13.
}
\]
\end{answerbox}

\end{workedexamplebox}

%%%%%%%%%%%%%%%%%%%%%%%%%%%%%%%%%%%%%%%%%%%%%%%%%%%%%%%%%%%%
\subsection{Integration Over General Domains}
\label{subsec:mc-general-domain-integration}
%%%%%%%%%%%%%%%%%%%%%%%%%%%%%%%%%%%%%%%%%%%%%%%%%%%%%%%%%%%%

Suppose \(D\subset\mathbb{R}^d\) has finite volume

\[
|D|.
\]

If

\[
X\sim\operatorname{Uniform}(D),
\]

then

\[
\boxed{
\int_D
f(x)\,dx
=
|D|
\mathbb{E}[f(X)].
}
\]

Thus,

\[
\boxed{
\widehat{I}_N
=
\frac{
|D|
}{
N
}
\sum_{i=1}^{N}
f(X_i).
}
\]

%%%%%%%%%%%%%%%%%%%%%%%%%%%%%%%%%%%%%%%%%%%%%%%%%%%%%%%%%%%%
\subsection{High-Dimensional Integration}
\label{subsec:mc-high-dimensional-integration}
%%%%%%%%%%%%%%%%%%%%%%%%%%%%%%%%%%%%%%%%%%%%%%%%%%%%%%%%%%%%

Traditional grid-based quadrature suffers from the curse of
dimensionality.

If \(m\) points are used in each of \(d\) dimensions, then the grid
contains

\[
\boxed{
m^d
}
\]

points.

Monte Carlo convergence remains approximately

\[
\boxed{
O(N^{-1/2})
}
\]

regardless of dimension.

The constant may still depend strongly on the integrand and dimension, but
the exponent does not deteriorate directly with \(d\).

%%%%%%%%%%%%%%%%%%%%%%%%%%%%%%%%%%%%%%%%%%%%%%%%%%%%%%%%%%%%
\subsection{Estimating Area and Volume}
\label{subsec:mc-area-volume}
%%%%%%%%%%%%%%%%%%%%%%%%%%%%%%%%%%%%%%%%%%%%%%%%%%%%%%%%%%%%

Suppose region \(A\) lies inside an easy-to-sample region \(B\).

If

\[
X\sim\operatorname{Uniform}(B),
\]

then

\[
\boxed{
P(X\in A)
=
\frac{
|A|
}{
|B|
}.
}
\]

Therefore,

\[
\boxed{
|A|
=
|B|
P(X\in A).
}
\]

Estimate the probability using the fraction of sampled points falling in
\(A\).

%%%%%%%%%%%%%%%%%%%%%%%%%%%%%%%%%%%%%%%%%%%%%%%%%%%%%%%%%%%%
\subsection{Estimating \(\pi\)}
\label{subsec:mc-estimating-pi}
%%%%%%%%%%%%%%%%%%%%%%%%%%%%%%%%%%%%%%%%%%%%%%%%%%%%%%%%%%%%

Sample uniformly from the square

\[
[-1,1]^2.
\]

The unit disk has area

\[
\pi,
\]

while the square has area

\[
4.
\]

Therefore,

\[
P(X^2+Y^2\leq1)
=
\frac{\pi}{4}.
\]

Hence,

\[
\boxed{
\pi
=
4
P(X^2+Y^2\leq1).
}
\]

The Monte Carlo estimator is

\[
\boxed{
\widehat{\pi}
=
4
\frac{
\#\{
i:
X_i^2+Y_i^2\leq1
\}
}{
N
}.
}
\]

%%%%%%%%%%%%%%%%%%%%%%%%%%%%%%%%%%%%%%%%%%%%%%%%%%%%%%%%%%%%
\subsection{Indicator Variables}
\label{subsec:mc-indicator-variables}
%%%%%%%%%%%%%%%%%%%%%%%%%%%%%%%%%%%%%%%%%%%%%%%%%%%%%%%%%%%%

For event \(A\), define

\[
\boxed{
\mathbf{1}_A(X)
=
\begin{cases}
1,
&
X\in A,
\\
0,
&
X\notin A.
\end{cases}
}
\]

Then

\[
\boxed{
\mathbb{E}
[
\mathbf{1}_A(X)
]
=
P(A).
}
\]

Thus event probabilities can always be written as expectations.

%%%%%%%%%%%%%%%%%%%%%%%%%%%%%%%%%%%%%%%%%%%%%%%%%%%%%%%%%%%%
\subsection{Probability Estimation}
\label{subsec:mc-probability-estimation}
%%%%%%%%%%%%%%%%%%%%%%%%%%%%%%%%%%%%%%%%%%%%%%%%%%%%%%%%%%%%

If

\[
X_i
\]

are independent samples, estimate

\[
p=P(A)
\]

using

\[
\boxed{
\widehat{p}
=
\frac1N
\sum_{i=1}^{N}
\mathbf{1}_A(X_i).
}
\]

Since the indicator is Bernoulli,

\[
\boxed{
\operatorname{Var}
(
\widehat{p}
)
=
\frac{
p(1-p)
}{
N
}.
}
\]

%%%%%%%%%%%%%%%%%%%%%%%%%%%%%%%%%%%%%%%%%%%%%%%%%%%%%%%%%%%%
\subsection{Random Number Generation}
\label{subsec:random-number-generation}
%%%%%%%%%%%%%%%%%%%%%%%%%%%%%%%%%%%%%%%%%%%%%%%%%%%%%%%%%%%%

Computer simulations usually rely on pseudorandom numbers.

A pseudorandom number generator produces a deterministic sequence designed
to imitate important statistical properties of random samples.

Given an initial internal state or seed, the sequence is reproducible.

%%%%%%%%%%%%%%%%%%%%%%%%%%%%%%%%%%%%%%%%%%%%%%%%%%%%%%%%%%%%
\subsection{Random Seeds}
\label{subsec:random-seeds}
%%%%%%%%%%%%%%%%%%%%%%%%%%%%%%%%%%%%%%%%%%%%%%%%%%%%%%%%%%%%

A random seed initializes a pseudorandom generator.

Using the same generator and seed generally reproduces the same sequence,
subject to implementation details.

This is useful for:

\[
\boxed{
\text{debugging},
}
\]

\[
\boxed{
\text{reproducibility},
}
\]

and

\[
\boxed{
\text{method comparison}.
}
\]

%%%%%%%%%%%%%%%%%%%%%%%%%%%%%%%%%%%%%%%%%%%%%%%%%%%%%%%%%%%%
\subsection{Uniform Random Variables}
\label{subsec:uniform-random-generation}
%%%%%%%%%%%%%%%%%%%%%%%%%%%%%%%%%%%%%%%%%%%%%%%%%%%%%%%%%%%%

Many simulation algorithms begin with

\[
\boxed{
U\sim\operatorname{Uniform}(0,1).
}
\]

More complicated distributions are then generated by transforming one or
more uniform random variables.

Thus uniform pseudorandom generation is a fundamental computational
primitive.

%%%%%%%%%%%%%%%%%%%%%%%%%%%%%%%%%%%%%%%%%%%%%%%%%%%%%%%%%%%%
\subsection{Inverse-Transform Sampling}
\label{subsec:inverse-transform-sampling}
%%%%%%%%%%%%%%%%%%%%%%%%%%%%%%%%%%%%%%%%%%%%%%%%%%%%%%%%%%%%

Suppose a continuous distribution has cumulative distribution function

\[
F(x).
\]

If

\[
U\sim\operatorname{Uniform}(0,1),
\]

then

\[
\boxed{
X
=
F^{-1}(U)
}
\]

has cumulative distribution function \(F\), assuming the inverse is
interpreted appropriately.

%%%%%%%%%%%%%%%%%%%%%%%%%%%%%%%%%%%%%%%%%%%%%%%%%%%%%%%%%%%%
\subsection{Worked Example: Exponential Simulation}
\label{subsec:worked-exponential-simulation}
%%%%%%%%%%%%%%%%%%%%%%%%%%%%%%%%%%%%%%%%%%%%%%%%%%%%%%%%%%%%

\begin{workedexamplebox}{Inverse Transform for an Exponential Variable}

Let

\[
X\sim\operatorname{Exponential}(\lambda).
\]

Its CDF is

\[
F(x)
=
1-e^{-\lambda x},
\qquad
x\geq0.
\]

Set

\[
U=F(X).
\]

Then

\[
U
=
1-e^{-\lambda X}.
\]

Thus,

\[
e^{-\lambda X}
=
1-U.
\]

Taking logarithms,

\[
X
=
-\frac1\lambda
\ln(1-U).
\]

Since \(1-U\) is also uniform on \((0,1)\), one may equivalently use

\[
\begin{answerbox}
\[
\boxed{
X
=
-\frac1\lambda
\ln U.
}
\]
\end{answerbox}

\end{workedexamplebox}

%%%%%%%%%%%%%%%%%%%%%%%%%%%%%%%%%%%%%%%%%%%%%%%%%%%%%%%%%%%%
\subsection{Transformation Method}
\label{subsec:mc-transformation-method}
%%%%%%%%%%%%%%%%%%%%%%%%%%%%%%%%%%%%%%%%%%%%%%%%%%%%%%%%%%%%

If a random variable \(X\) has an easy-to-simulate distribution and

\[
Y=g(X),
\]

then samples of \(Y\) can be obtained by evaluating

\[
\boxed{
Y_i=g(X_i).
}
\]

Transformation methods are useful when the desired distribution can be
constructed from simpler distributions.

%%%%%%%%%%%%%%%%%%%%%%%%%%%%%%%%%%%%%%%%%%%%%%%%%%%%%%%%%%%%
\subsection{Acceptance--Rejection Sampling}
\label{subsec:acceptance-rejection}
%%%%%%%%%%%%%%%%%%%%%%%%%%%%%%%%%%%%%%%%%%%%%%%%%%%%%%%%%%%%

Suppose target density \(f\) is difficult to sample from directly.

Choose an easy proposal density \(g\) and constant \(M\) such that

\[
\boxed{
f(x)
\leq
Mg(x)
}
\]

for all relevant \(x\).

Then:

\begin{enumerate}
    \item sample \(Y\sim g\);
    \item sample \(U\sim\operatorname{Uniform}(0,1)\);
    \item accept \(Y\) if
    \[
    U
    \leq
    \frac{
    f(Y)
    }{
    Mg(Y)
    }.
    \]
\end{enumerate}

Accepted values follow the target distribution.

%%%%%%%%%%%%%%%%%%%%%%%%%%%%%%%%%%%%%%%%%%%%%%%%%%%%%%%%%%%%
\subsection{Acceptance Probability}
\label{subsec:acceptance-probability}
%%%%%%%%%%%%%%%%%%%%%%%%%%%%%%%%%%%%%%%%%%%%%%%%%%%%%%%%%%%%

If \(f\) is normalized and

\[
f(x)\leq Mg(x),
\]

then the average acceptance probability is

\[
\boxed{
\frac1M.
}
\]

Thus a large \(M\) produces many rejected samples.

A good proposal should closely match the target.

%%%%%%%%%%%%%%%%%%%%%%%%%%%%%%%%%%%%%%%%%%%%%%%%%%%%%%%%%%%%
\subsection{Variance Reduction}
\label{subsec:variance-reduction}
%%%%%%%%%%%%%%%%%%%%%%%%%%%%%%%%%%%%%%%%%%%%%%%%%%%%%%%%%%%%

The standard Monte Carlo error decreases only as

\[
N^{-1/2}.
\]

Rather than increasing \(N\), one can sometimes reduce the variance of the
estimator.

Variance reduction seeks

\[
\boxed{
\text{same expectation}
+
\text{smaller variance}.
}
\]

Important techniques include:

\[
\boxed{
\text{importance sampling},
}
\]

\[
\boxed{
\text{control variates},
}
\]

\[
\boxed{
\text{antithetic variates},
}
\]

and

\[
\boxed{
\text{stratification}.
}
\]

%%%%%%%%%%%%%%%%%%%%%%%%%%%%%%%%%%%%%%%%%%%%%%%%%%%%%%%%%%%%
\subsection{Importance Sampling}
\label{subsec:importance-sampling}
%%%%%%%%%%%%%%%%%%%%%%%%%%%%%%%%%%%%%%%%%%%%%%%%%%%%%%%%%%%%

Suppose

\[
\mu
=
\int
g(x)
p(x)\,dx.
\]

Choose another density \(q(x)\) satisfying

\[
q(x)>0
\]

whenever

\[
g(x)p(x)\neq0.
\]

Then

\[
\mu
=
\int
g(x)
\frac{
p(x)
}{
q(x)
}
q(x)\,dx.
\]

Therefore,

\[
\boxed{
\mu
=
\mathbb{E}_q
\left[
g(X)
\frac{
p(X)
}{
q(X)
}
\right].
}
\]

%%%%%%%%%%%%%%%%%%%%%%%%%%%%%%%%%%%%%%%%%%%%%%%%%%%%%%%%%%%%
\subsection{Importance Weights}
\label{subsec:importance-weights}
%%%%%%%%%%%%%%%%%%%%%%%%%%%%%%%%%%%%%%%%%%%%%%%%%%%%%%%%%%%%

Define the importance weight

\[
\boxed{
w(x)
=
\frac{
p(x)
}{
q(x)
}.
}
\]

Then

\[
\boxed{
\widehat{\mu}_{IS}
=
\frac1N
\sum_{i=1}^{N}
g(X_i)
w(X_i),
\qquad
X_i\sim q.
}
\]

A good proposal places more probability in regions contributing strongly
to the expectation.

%%%%%%%%%%%%%%%%%%%%%%%%%%%%%%%%%%%%%%%%%%%%%%%%%%%%%%%%%%%%
\subsection{Importance-Sampling Variance}
\label{subsec:importance-sampling-variance}
%%%%%%%%%%%%%%%%%%%%%%%%%%%%%%%%%%%%%%%%%%%%%%%%%%%%%%%%%%%%

The variance is

\[
\boxed{
\operatorname{Var}_q
\left(
g(X)\frac{p(X)}{q(X)}
\right).
}
\]

A poor proposal can make this variance extremely large or even infinite.

Thus importance sampling is powerful but can perform much worse than
ordinary Monte Carlo if the proposal is badly chosen.

%%%%%%%%%%%%%%%%%%%%%%%%%%%%%%%%%%%%%%%%%%%%%%%%%%%%%%%%%%%%
\subsection{Rare Events}
\label{subsec:rare-event-monte-carlo}
%%%%%%%%%%%%%%%%%%%%%%%%%%%%%%%%%%%%%%%%%%%%%%%%%%%%%%%%%%%%

Suppose

\[
p=P(A)
\]

is very small.

The naive estimator is

\[
\widehat{p}
=
\frac1N
\sum_{i=1}^{N}
\mathbf{1}_A(X_i).
\]

Its relative standard error is approximately

\[
\boxed{
\frac{
\sqrt{p(1-p)/N}
}{
p
}
=
\sqrt{
\frac{
1-p
}{
Np
}
}.
}
\]

For very small \(p\), this is large unless

\[
Np
\]

is itself large.

%%%%%%%%%%%%%%%%%%%%%%%%%%%%%%%%%%%%%%%%%%%%%%%%%%%%%%%%%%%%
\subsection{Importance Sampling for Rare Events}
\label{subsec:importance-sampling-rare-events}
%%%%%%%%%%%%%%%%%%%%%%%%%%%%%%%%%%%%%%%%%%%%%%%%%%%%%%%%%%%%

Rare-event importance sampling changes the sampling distribution so the
rare event occurs much more frequently.

The likelihood ratio

\[
\boxed{
\frac{p(x)}{q(x)}
}
\]

then corrects for the altered sampling distribution.

A successful proposal substantially reduces estimator variance while
remaining unbiased under standard conditions.

%%%%%%%%%%%%%%%%%%%%%%%%%%%%%%%%%%%%%%%%%%%%%%%%%%%%%%%%%%%%
\subsection{Control Variates}
\label{subsec:control-variates}
%%%%%%%%%%%%%%%%%%%%%%%%%%%%%%%%%%%%%%%%%%%%%%%%%%%%%%%%%%%%

Suppose we want

\[
\mu_X
=
\mathbb{E}[X],
\]

and have another correlated variable \(Y\) with known expectation

\[
\mu_Y.
\]

Define

\[
\boxed{
Z
=
X
-
c
(
Y-\mu_Y
).
}
\]

Then

\[
\boxed{
\mathbb{E}[Z]
=
\mu_X.
}
\]

Thus \(Z\) is another unbiased estimator target.

%%%%%%%%%%%%%%%%%%%%%%%%%%%%%%%%%%%%%%%%%%%%%%%%%%%%%%%%%%%%
\subsection{Optimal Control-Variate Coefficient}
\label{subsec:optimal-control-variate}
%%%%%%%%%%%%%%%%%%%%%%%%%%%%%%%%%%%%%%%%%%%%%%%%%%%%%%%%%%%%

The variance is minimized by

\[
\boxed{
c^*
=
\frac{
\operatorname{Cov}(X,Y)
}{
\operatorname{Var}(Y)
}
}
\]

when

\[
\operatorname{Var}(Y)>0.
\]

Then

\[
\boxed{
\operatorname{Var}(Z)
=
\operatorname{Var}(X)
\left(
1-\rho_{XY}^2
\right),
}
\]

where \(\rho_{XY}\) is the correlation.

Strong correlation can produce major variance reduction.

%%%%%%%%%%%%%%%%%%%%%%%%%%%%%%%%%%%%%%%%%%%%%%%%%%%%%%%%%%%%
\subsection{Worked Example: Control Variate}
\label{subsec:worked-control-variate}
%%%%%%%%%%%%%%%%%%%%%%%%%%%%%%%%%%%%%%%%%%%%%%%%%%%%%%%%%%%%

\begin{workedexamplebox}{Using a Known Mean}

Suppose \(X\) is expensive to simulate and is positively correlated with
\(Y\), where

\[
\mathbb{E}[Y]=10.
\]

Choose coefficient \(c=2\).

Then use

\[
Z
=
X
-
2(Y-10).
\]

Because

\[
\mathbb{E}[Y-10]=0,
\]

we have

\[
\mathbb{E}[Z]
=
\mathbb{E}[X].
\]

\begin{answerbox}
\[
\boxed{
\widehat{\mu}_{CV}
=
\frac1N
\sum_{i=1}^{N}
\left[
X_i-2(Y_i-10)
\right].
}
\]
\end{answerbox}

If \(X\) and \(Y\) are strongly correlated, the estimator may have much
smaller variance than the ordinary sample mean of \(X\).

\end{workedexamplebox}

%%%%%%%%%%%%%%%%%%%%%%%%%%%%%%%%%%%%%%%%%%%%%%%%%%%%%%%%%%%%
\subsection{Antithetic Variates}
\label{subsec:antithetic-variates}
%%%%%%%%%%%%%%%%%%%%%%%%%%%%%%%%%%%%%%%%%%%%%%%%%%%%%%%%%%%%

Suppose

\[
U\sim\operatorname{Uniform}(0,1).
\]

Then

\[
1-U
\]

has the same distribution.

Instead of using only

\[
g(U),
\]

form the pair

\[
\boxed{
g(U),
\qquad
g(1-U).
}
\]

The antithetic estimator is

\[
\boxed{
\widehat{\mu}
=
\frac1N
\sum_{i=1}^{N}
\frac{
g(U_i)+g(1-U_i)
}{
2
}.
}
\]

%%%%%%%%%%%%%%%%%%%%%%%%%%%%%%%%%%%%%%%%%%%%%%%%%%%%%%%%%%%%
\subsection{Why Antithetic Variates Can Help}
\label{subsec:why-antithetic-helps}
%%%%%%%%%%%%%%%%%%%%%%%%%%%%%%%%%%%%%%%%%%%%%%%%%%%%%%%%%%%%

If

\[
g(U)
\]

and

\[
g(1-U)
\]

are negatively correlated, averaging them reduces variance.

The method is especially natural for monotone transformations of uniform
random variables.

Negative correlation is the key mechanism.

%%%%%%%%%%%%%%%%%%%%%%%%%%%%%%%%%%%%%%%%%%%%%%%%%%%%%%%%%%%%
\subsection{Stratified Sampling}
\label{subsec:stratified-sampling}
%%%%%%%%%%%%%%%%%%%%%%%%%%%%%%%%%%%%%%%%%%%%%%%%%%%%%%%%%%%%

Divide the sample space into disjoint strata

\[
A_1,\ldots,A_K.
\]

Then

\[
\boxed{
\mathbb{E}[g(X)]
=
\sum_{k=1}^{K}
P(A_k)
\mathbb{E}
[
g(X)
\mid
A_k
].
}
\]

Estimate each conditional expectation separately.

This ensures that important portions of the sample space receive
representation.

%%%%%%%%%%%%%%%%%%%%%%%%%%%%%%%%%%%%%%%%%%%%%%%%%%%%%%%%%%%%
\subsection{Stratified Estimator}
\label{subsec:stratified-estimator}
%%%%%%%%%%%%%%%%%%%%%%%%%%%%%%%%%%%%%%%%%%%%%%%%%%%%%%%%%%%%

Let

\[
p_k=P(A_k).
\]

If \(N_k\) samples are drawn within stratum \(k\), then

\[
\boxed{
\widehat{\mu}_{strat}
=
\sum_{k=1}^{K}
p_k
\widehat{\mu}_k,
}
\]

where

\[
\boxed{
\widehat{\mu}_k
=
\frac1{N_k}
\sum_{i=1}^{N_k}
g(X_{ki}).
}
\]

The allocation \(N_k\) affects variance.

%%%%%%%%%%%%%%%%%%%%%%%%%%%%%%%%%%%%%%%%%%%%%%%%%%%%%%%%%%%%
\subsection{Latin Hypercube Sampling Preview}
\label{subsec:latin-hypercube}
%%%%%%%%%%%%%%%%%%%%%%%%%%%%%%%%%%%%%%%%%%%%%%%%%%%%%%%%%%%%

Latin hypercube sampling divides each input dimension into intervals of
equal probability.

Samples are arranged so that each interval is represented exactly once in
each dimension.

This provides more uniform coverage than ordinary independent random
sampling in many uncertainty-quantification problems.

%%%%%%%%%%%%%%%%%%%%%%%%%%%%%%%%%%%%%%%%%%%%%%%%%%%%%%%%%%%%
\subsection{Common Random Numbers}
\label{subsec:common-random-numbers}
%%%%%%%%%%%%%%%%%%%%%%%%%%%%%%%%%%%%%%%%%%%%%%%%%%%%%%%%%%%%

Suppose two systems are compared using simulated performance

\[
Y_A
\]

and

\[
Y_B.
\]

Use the same underlying random inputs when simulating both systems.

Then estimate

\[
\boxed{
\mathbb{E}
[
Y_A-Y_B
].
}
\]

If the outputs respond similarly to common randomness, their difference
may have substantially reduced variance.

%%%%%%%%%%%%%%%%%%%%%%%%%%%%%%%%%%%%%%%%%%%%%%%%%%%%%%%%%%%%
\subsection{Monte Carlo for Optimization}
\label{subsec:monte-carlo-optimization}
%%%%%%%%%%%%%%%%%%%%%%%%%%%%%%%%%%%%%%%%%%%%%%%%%%%%%%%%%%%%

Suppose an objective is

\[
\boxed{
F(x)
=
\mathbb{E}
[
f(x,\xi)
].
}
\]

Monte Carlo replaces the expectation with a sample average:

\[
\boxed{
\widehat{F}_N(x)
=
\frac1N
\sum_{i=1}^{N}
f(x,\xi_i).
}
\]

Then solve

\[
\boxed{
\min_x
\widehat{F}_N(x).
}
\]

This is a Sample Average Approximation approach.

%%%%%%%%%%%%%%%%%%%%%%%%%%%%%%%%%%%%%%%%%%%%%%%%%%%%%%%%%%%%
\subsection{Monte Carlo Gradient Estimation}
\label{subsec:monte-carlo-gradients}
%%%%%%%%%%%%%%%%%%%%%%%%%%%%%%%%%%%%%%%%%%%%%%%%%%%%%%%%%%%%

If differentiation and expectation can be interchanged under appropriate
regularity conditions,

\[
\boxed{
\nabla F(x)
=
\mathbb{E}
[
\nabla_x f(x,\xi)
].
}
\]

Therefore,

\[
\boxed{
\widehat{\nabla F}_N(x)
=
\frac1N
\sum_{i=1}^{N}
\nabla_x f(x,\xi_i).
}
\]

This connects Monte Carlo directly with stochastic gradient methods.

%%%%%%%%%%%%%%%%%%%%%%%%%%%%%%%%%%%%%%%%%%%%%%%%%%%%%%%%%%%%
\subsection{Markov Chain Monte Carlo}
\label{subsec:markov-chain-monte-carlo}
%%%%%%%%%%%%%%%%%%%%%%%%%%%%%%%%%%%%%%%%%%%%%%%%%%%%%%%%%%%%

Ordinary Monte Carlo often assumes independent samples.

Sometimes direct independent sampling from a complicated target density

\[
\pi(x)
\]

is difficult.

Markov Chain Monte Carlo, abbreviated MCMC, constructs a Markov chain

\[
X_0,X_1,X_2,\ldots
\]

whose stationary distribution is the desired target \(\pi\).

%%%%%%%%%%%%%%%%%%%%%%%%%%%%%%%%%%%%%%%%%%%%%%%%%%%%%%%%%%%%
\subsection{Stationary Distribution}
\label{subsec:mcmc-stationary-distribution}
%%%%%%%%%%%%%%%%%%%%%%%%%%%%%%%%%%%%%%%%%%%%%%%%%%%%%%%%%%%%

A distribution \(\pi\) is stationary for transition kernel \(P\) if

\[
\boxed{
\pi P=\pi
}
\]

in discrete notation.

If

\[
X_n\sim\pi,
\]

then

\[
X_{n+1}\sim\pi.
\]

A properly designed MCMC algorithm attempts to approach this stationary
distribution as the chain evolves.

%%%%%%%%%%%%%%%%%%%%%%%%%%%%%%%%%%%%%%%%%%%%%%%%%%%%%%%%%%%%
\subsection{MCMC Estimator}
\label{subsec:mcmc-estimator}
%%%%%%%%%%%%%%%%%%%%%%%%%%%%%%%%%%%%%%%%%%%%%%%%%%%%%%%%%%%%

After generating

\[
X_1,\ldots,X_N,
\]

estimate

\[
\mathbb{E}_\pi[g(X)]
\]

using

\[
\boxed{
\widehat{\mu}
=
\frac1N
\sum_{i=1}^{N}
g(X_i).
}
\]

Unlike ordinary independent Monte Carlo, the samples are generally
correlated.

Appropriate ergodic conditions replace the ordinary independent-sample
Law of Large Numbers.

%%%%%%%%%%%%%%%%%%%%%%%%%%%%%%%%%%%%%%%%%%%%%%%%%%%%%%%%%%%%
\subsection{Metropolis--Hastings Algorithm}
\label{subsec:metropolis-hastings}
%%%%%%%%%%%%%%%%%%%%%%%%%%%%%%%%%%%%%%%%%%%%%%%%%%%%%%%%%%%%

Suppose the current state is

\[
X_n=x.
\]

Choose a proposal distribution

\[
\boxed{
q(y\mid x).
}
\]

Generate candidate

\[
Y\sim q(\cdot\mid x).
\]

Accept the candidate with probability

\[
\boxed{
\alpha(x,y)
=
\min
\left\{
1,
\frac{
\pi(y)q(x\mid y)
}{
\pi(x)q(y\mid x)
}
\right\}.
}
\]

If rejected, remain at \(x\).

%%%%%%%%%%%%%%%%%%%%%%%%%%%%%%%%%%%%%%%%%%%%%%%%%%%%%%%%%%%%
\subsection{Metropolis--Hastings Update}
\label{subsec:metropolis-hastings-update}
%%%%%%%%%%%%%%%%%%%%%%%%%%%%%%%%%%%%%%%%%%%%%%%%%%%%%%%%%%%%

Generate

\[
U\sim\operatorname{Uniform}(0,1).
\]

Then

\[
\boxed{
X_{n+1}
=
\begin{cases}
Y,
&
U\leq\alpha(x,Y),
\\
x,
&
U>\alpha(x,Y).
\end{cases}
}
\]

Repeated application produces a dependent sequence of states.

%%%%%%%%%%%%%%%%%%%%%%%%%%%%%%%%%%%%%%%%%%%%%%%%%%%%%%%%%%%%
\subsection{Symmetric Proposal}
\label{subsec:symmetric-mh-proposal}
%%%%%%%%%%%%%%%%%%%%%%%%%%%%%%%%%%%%%%%%%%%%%%%%%%%%%%%%%%%%

If

\[
q(y\mid x)
=
q(x\mid y),
\]

then the proposal terms cancel.

The acceptance probability becomes

\[
\boxed{
\alpha(x,y)
=
\min
\left\{
1,
\frac{
\pi(y)
}{
\pi(x)
}
\right\}.
}
\]

This is the classical Metropolis form.

%%%%%%%%%%%%%%%%%%%%%%%%%%%%%%%%%%%%%%%%%%%%%%%%%%%%%%%%%%%%
\subsection{Unnormalized Target Densities}
\label{subsec:unnormalized-mcmc-target}
%%%%%%%%%%%%%%%%%%%%%%%%%%%%%%%%%%%%%%%%%%%%%%%%%%%%%%%%%%%%

Suppose

\[
\pi(x)
=
\frac{
\widetilde{\pi}(x)
}{
Z
},
\]

where normalizing constant \(Z\) is unknown.

In the Metropolis--Hastings ratio,

\[
Z
\]

cancels.

Thus MCMC can sample from many distributions known only up to a
multiplicative constant.

This is especially important in Bayesian inference.

%%%%%%%%%%%%%%%%%%%%%%%%%%%%%%%%%%%%%%%%%%%%%%%%%%%%%%%%%%%%
\subsection{Worked Example: Metropolis Acceptance}
\label{subsec:worked-metropolis-acceptance}
%%%%%%%%%%%%%%%%%%%%%%%%%%%%%%%%%%%%%%%%%%%%%%%%%%%%%%%%%%%%

\begin{workedexamplebox}{Computing an Acceptance Probability}

Suppose a symmetric proposal is used and

\[
\pi(x)=0.20,
\]

while the proposed state's target density is

\[
\pi(y)=0.10.
\]

Then

\[
\alpha(x,y)
=
\min
\left\{
1,
\frac{0.10}{0.20}
\right\}.
\]

Therefore,

\[
\begin{answerbox}
\[
\boxed{
\alpha(x,y)=0.5.
}
\]
\end{answerbox}

The proposed move is accepted with probability \(0.5\).

\end{workedexamplebox}

%%%%%%%%%%%%%%%%%%%%%%%%%%%%%%%%%%%%%%%%%%%%%%%%%%%%%%%%%%%%
\subsection{Detailed Balance}
\label{subsec:detailed-balance}
%%%%%%%%%%%%%%%%%%%%%%%%%%%%%%%%%%%%%%%%%%%%%%%%%%%%%%%%%%%%

A sufficient condition for stationarity is detailed balance:

\[
\boxed{
\pi(x)
P(x,dy)
=
\pi(y)
P(y,dx).
}
\]

This expresses reversible probability flow between states.

Metropolis--Hastings is designed so that the target distribution satisfies
this condition under standard assumptions.

%%%%%%%%%%%%%%%%%%%%%%%%%%%%%%%%%%%%%%%%%%%%%%%%%%%%%%%%%%%%
\subsection{Gibbs Sampling}
\label{subsec:gibbs-sampling}
%%%%%%%%%%%%%%%%%%%%%%%%%%%%%%%%%%%%%%%%%%%%%%%%%%%%%%%%%%%%

Suppose

\[
X
=
(X_1,\ldots,X_d)
\]

has a complicated joint distribution, but conditional distributions are
easy to sample.

Gibbs sampling updates components sequentially from

\[
\boxed{
X_j
\sim
\pi
(
x_j
\mid
x_{-j}
).
}
\]

A full sweep updates each component using its current conditional
distribution.

%%%%%%%%%%%%%%%%%%%%%%%%%%%%%%%%%%%%%%%%%%%%%%%%%%%%%%%%%%%%
\subsection{Gibbs Sampling as MCMC}
\label{subsec:gibbs-as-mcmc}
%%%%%%%%%%%%%%%%%%%%%%%%%%%%%%%%%%%%%%%%%%%%%%%%%%%%%%%%%%%%

Gibbs sampling is a special MCMC method.

Each conditional update leaves the target distribution invariant under
appropriate conditions.

Unlike general Metropolis--Hastings, Gibbs updates are accepted
automatically.

%%%%%%%%%%%%%%%%%%%%%%%%%%%%%%%%%%%%%%%%%%%%%%%%%%%%%%%%%%%%
\subsection{Burn-In}
\label{subsec:mcmc-burn-in}
%%%%%%%%%%%%%%%%%%%%%%%%%%%%%%%%%%%%%%%%%%%%%%%%%%%%%%%%%%%%

The initial MCMC state

\[
X_0
\]

may be far from typical regions of the target distribution.

Early iterations are sometimes discarded as burn-in.

Burn-in does not by itself guarantee convergence.

The more important issue is whether the chain has reached representative
target behavior.

%%%%%%%%%%%%%%%%%%%%%%%%%%%%%%%%%%%%%%%%%%%%%%%%%%%%%%%%%%%%
\subsection{MCMC Autocorrelation}
\label{subsec:mcmc-autocorrelation}
%%%%%%%%%%%%%%%%%%%%%%%%%%%%%%%%%%%%%%%%%%%%%%%%%%%%%%%%%%%%

MCMC samples are correlated.

For a scalar quantity

\[
Y_n=g(X_n),
\]

define lag-\(k\) autocorrelation

\[
\boxed{
\rho_k
=
\operatorname{Corr}
(
Y_n,Y_{n+k}
).
}
\]

Strong positive autocorrelation means the chain produces less information
than the same number of independent samples.

%%%%%%%%%%%%%%%%%%%%%%%%%%%%%%%%%%%%%%%%%%%%%%%%%%%%%%%%%%%%
\subsection{Integrated Autocorrelation Time}
\label{subsec:integrated-autocorrelation-time}
%%%%%%%%%%%%%%%%%%%%%%%%%%%%%%%%%%%%%%%%%%%%%%%%%%%%%%%%%%%%

A common measure is

\[
\boxed{
\tau_{\mathrm{int}}
=
1
+
2
\sum_{k=1}^{\infty}
\rho_k,
}
\]

when the sum is well defined.

The variance of the MCMC average is inflated approximately by this factor
relative to independent sampling.

%%%%%%%%%%%%%%%%%%%%%%%%%%%%%%%%%%%%%%%%%%%%%%%%%%%%%%%%%%%%
\subsection{Effective Sample Size}
\label{subsec:effective-sample-size}
%%%%%%%%%%%%%%%%%%%%%%%%%%%%%%%%%%%%%%%%%%%%%%%%%%%%%%%%%%%%

An approximate effective sample size is

\[
\boxed{
N_{\mathrm{eff}}
\approx
\frac{
N
}{
\tau_{\mathrm{int}}
}.
}
\]

If

\[
N=10000
\]

but

\[
\tau_{\mathrm{int}}=20,
\]

then

\[
N_{\mathrm{eff}}
\approx500.
\]

Thus chain length alone does not measure information content.

%%%%%%%%%%%%%%%%%%%%%%%%%%%%%%%%%%%%%%%%%%%%%%%%%%%%%%%%%%%%
\subsection{MCMC Mixing}
\label{subsec:mcmc-mixing}
%%%%%%%%%%%%%%%%%%%%%%%%%%%%%%%%%%%%%%%%%%%%%%%%%%%%%%%%%%%%

A chain mixes well if it explores the target distribution efficiently.

Poor mixing may appear as:

\[
\boxed{
\text{long periods in one region},
}
\]

\[
\boxed{
\text{high autocorrelation},
}
\]

or

\[
\boxed{
\text{failure to move between important modes}.
}
\]

Proposal design strongly affects mixing.

%%%%%%%%%%%%%%%%%%%%%%%%%%%%%%%%%%%%%%%%%%%%%%%%%%%%%%%%%%%%
\subsection{Acceptance Rate}
\label{subsec:mcmc-acceptance-rate}
%%%%%%%%%%%%%%%%%%%%%%%%%%%%%%%%%%%%%%%%%%%%%%%%%%%%%%%%%%%%

If proposals are extremely small, acceptance may be high but exploration
slow.

If proposals are extremely large, many moves may be rejected.

Thus acceptance rate alone is not a sufficient diagnostic.

A good algorithm balances movement and acceptance.

%%%%%%%%%%%%%%%%%%%%%%%%%%%%%%%%%%%%%%%%%%%%%%%%%%%%%%%%%%%%
\subsection{Multiple Chains}
\label{subsec:mcmc-multiple-chains}
%%%%%%%%%%%%%%%%%%%%%%%%%%%%%%%%%%%%%%%%%%%%%%%%%%%%%%%%%%%%

Running several chains from different starting points can reveal whether
the simulation explores similar regions of the target distribution.

Agreement across chains provides useful evidence of convergence.

Disagreement may indicate poor mixing, multimodality, or inadequate run
length.

%%%%%%%%%%%%%%%%%%%%%%%%%%%%%%%%%%%%%%%%%%%%%%%%%%%%%%%%%%%%
\subsection{Trace Plots}
\label{subsec:mcmc-trace-plots}
%%%%%%%%%%%%%%%%%%%%%%%%%%%%%%%%%%%%%%%%%%%%%%%%%%%%%%%%%%%%

A trace plot displays

\[
g(X_n)
\]

against iteration number \(n\).

A well-mixing chain often shows repeated movement through its typical
range without long systematic drifts.

Trace plots are useful but should not be the only diagnostic.

%%%%%%%%%%%%%%%%%%%%%%%%%%%%%%%%%%%%%%%%%%%%%%%%%%%%%%%%%%%%
\subsection{Independent Monte Carlo Versus MCMC}
\label{subsec:independent-mc-vs-mcmc}
%%%%%%%%%%%%%%%%%%%%%%%%%%%%%%%%%%%%%%%%%%%%%%%%%%%%%%%%%%%%

Independent Monte Carlo uses

\[
\boxed{
X_i
\overset{\mathrm{iid}}{\sim}
\pi.
}
\]

MCMC uses

\[
\boxed{
X_{i+1}
\sim
P(\cdot\mid X_i),
}
\]

with stationary distribution \(\pi\).

Independent sampling simplifies variance analysis.

MCMC allows access to difficult target distributions at the price of
dependence and more complex diagnostics.

%%%%%%%%%%%%%%%%%%%%%%%%%%%%%%%%%%%%%%%%%%%%%%%%%%%%%%%%%%%%
\subsection{Sequential Monte Carlo Preview}
\label{subsec:sequential-monte-carlo}
%%%%%%%%%%%%%%%%%%%%%%%%%%%%%%%%%%%%%%%%%%%%%%%%%%%%%%%%%%%%

Sequential Monte Carlo represents a distribution using a collection of
weighted particles:

\[
\boxed{
\{
X_i,w_i
\}_{i=1}^{N}.
}
\]

As the target distribution changes, particles are:

\[
\boxed{
\text{propagated},
}
\]

\[
\boxed{
\text{reweighted},
}
\]

and sometimes

\[
\boxed{
\text{resampled}.
}
\]

These methods are useful in filtering and dynamic Bayesian models.

%%%%%%%%%%%%%%%%%%%%%%%%%%%%%%%%%%%%%%%%%%%%%%%%%%%%%%%%%%%%
\subsection{Particle Degeneracy}
\label{subsec:particle-degeneracy}
%%%%%%%%%%%%%%%%%%%%%%%%%%%%%%%%%%%%%%%%%%%%%%%%%%%%%%%%%%%%

After repeated weighting, a small number of particles may contain nearly
all the total weight.

This is called weight degeneracy.

Resampling removes low-weight particles and replicates high-weight
particles.

However, excessive resampling can reduce diversity.

%%%%%%%%%%%%%%%%%%%%%%%%%%%%%%%%%%%%%%%%%%%%%%%%%%%%%%%%%%%%
\subsection{Bootstrap}
\label{subsec:bootstrap-monte-carlo}
%%%%%%%%%%%%%%%%%%%%%%%%%%%%%%%%%%%%%%%%%%%%%%%%%%%%%%%%%%%%

The bootstrap approximates the sampling distribution of a statistic by
resampling observed data.

Given observations

\[
X_1,\ldots,X_n,
\]

construct a bootstrap sample

\[
X_1^*,\ldots,X_n^*
\]

by sampling with replacement from the observed values.

Compute the statistic

\[
\boxed{
T^*
=
T(X_1^*,\ldots,X_n^*).
}
\]

Repeat many times.

%%%%%%%%%%%%%%%%%%%%%%%%%%%%%%%%%%%%%%%%%%%%%%%%%%%%%%%%%%%%
\subsection{Bootstrap Standard Error}
\label{subsec:bootstrap-standard-error}
%%%%%%%%%%%%%%%%%%%%%%%%%%%%%%%%%%%%%%%%%%%%%%%%%%%%%%%%%%%%

Suppose bootstrap statistics are

\[
T_1^*,\ldots,T_B^*.
\]

Their sample standard deviation estimates the standard error of the
original statistic:

\[
\boxed{
\widehat{\operatorname{SE}}_{\mathrm{boot}}
=
\sqrt{
\frac1{B-1}
\sum_{b=1}^{B}
\left(
T_b^*
-
\overline{T}^*
\right)^2
}.
}
\]

%%%%%%%%%%%%%%%%%%%%%%%%%%%%%%%%%%%%%%%%%%%%%%%%%%%%%%%%%%%%
\subsection{Parametric and Nonparametric Bootstrap}
\label{subsec:parametric-nonparametric-bootstrap}
%%%%%%%%%%%%%%%%%%%%%%%%%%%%%%%%%%%%%%%%%%%%%%%%%%%%%%%%%%%%

The nonparametric bootstrap resamples directly from the empirical data.

The parametric bootstrap fits a probability model and simulates new data
from the fitted distribution.

Thus:

\[
\boxed{
\text{nonparametric}
\rightarrow
\text{resample observations},
}
\]

while

\[
\boxed{
\text{parametric}
\rightarrow
\text{simulate fitted model}.
}
\]

%%%%%%%%%%%%%%%%%%%%%%%%%%%%%%%%%%%%%%%%%%%%%%%%%%%%%%%%%%%%
\subsection{Quasi-Monte Carlo}
\label{subsec:quasi-monte-carlo}
%%%%%%%%%%%%%%%%%%%%%%%%%%%%%%%%%%%%%%%%%%%%%%%%%%%%%%%%%%%%

Quasi-Monte Carlo replaces pseudorandom independent points with
deterministic low-discrepancy point sets.

The goal is to cover

\[
[0,1]^d
\]

more uniformly than random sampling.

A quasi-Monte Carlo integral has form

\[
\boxed{
\widehat{I}_N
=
\frac1N
\sum_{i=1}^{N}
f(x_i),
}
\]

where the \(x_i\) are carefully constructed deterministic points.

%%%%%%%%%%%%%%%%%%%%%%%%%%%%%%%%%%%%%%%%%%%%%%%%%%%%%%%%%%%%
\subsection{Discrepancy}
\label{subsec:quasi-monte-carlo-discrepancy}
%%%%%%%%%%%%%%%%%%%%%%%%%%%%%%%%%%%%%%%%%%%%%%%%%%%%%%%%%%%%

Discrepancy measures how uniformly a finite set of points covers a domain.

Low-discrepancy sequences avoid the random clustering and empty gaps that
naturally occur in independent Monte Carlo sampling.

Examples include:

\[
\boxed{
\text{Halton-type sequences},
}
\]

\[
\boxed{
\text{Sobol-type sequences},
}
\]

and other digital nets.

%%%%%%%%%%%%%%%%%%%%%%%%%%%%%%%%%%%%%%%%%%%%%%%%%%%%%%%%%%%%
\subsection{Monte Carlo Versus Quasi-Monte Carlo}
\label{subsec:mc-vs-qmc}
%%%%%%%%%%%%%%%%%%%%%%%%%%%%%%%%%%%%%%%%%%%%%%%%%%%%%%%%%%%%

Ordinary Monte Carlo is randomized and has a familiar probabilistic error
scale of

\[
O(N^{-1/2}).
\]

Quasi-Monte Carlo is deterministic and can converge faster for sufficiently
regular low-effective-dimensional integrands.

However, its error theory is different and depends on discrepancy and
integrand variation.

%%%%%%%%%%%%%%%%%%%%%%%%%%%%%%%%%%%%%%%%%%%%%%%%%%%%%%%%%%%%
\subsection{Randomized Quasi-Monte Carlo}
\label{subsec:randomized-qmc}
%%%%%%%%%%%%%%%%%%%%%%%%%%%%%%%%%%%%%%%%%%%%%%%%%%%%%%%%%%%%

Low-discrepancy point sets can be randomized while preserving their
uniformity properties.

This allows:

\[
\boxed{
\text{replicated estimates},
}
\]

and

\[
\boxed{
\text{empirical error assessment}.
}
\]

Randomized quasi-Monte Carlo combines deterministic space-filling
structure with probabilistic uncertainty estimation.

%%%%%%%%%%%%%%%%%%%%%%%%%%%%%%%%%%%%%%%%%%%%%%%%%%%%%%%%%%%%
\subsection{Monte Carlo for Stochastic Differential Equations}
\label{subsec:mc-sdes}
%%%%%%%%%%%%%%%%%%%%%%%%%%%%%%%%%%%%%%%%%%%%%%%%%%%%%%%%%%%%

Suppose

\[
\boxed{
dX_t
=
a(X_t,t)\,dt
+
b(X_t,t)\,dW_t.
}
\]

If the quantity of interest is

\[
\boxed{
\mathbb{E}
[
g(X_T)
],
}
\]

simulate many approximate sample paths and average:

\[
\boxed{
\widehat{\mu}_N
=
\frac1N
\sum_{i=1}^{N}
g
(
X_T^{(i)}
).
}
\]

%%%%%%%%%%%%%%%%%%%%%%%%%%%%%%%%%%%%%%%%%%%%%%%%%%%%%%%%%%%%
\subsection{Euler--Maruyama Method}
\label{subsec:euler-maruyama}
%%%%%%%%%%%%%%%%%%%%%%%%%%%%%%%%%%%%%%%%%%%%%%%%%%%%%%%%%%%%

A basic numerical method for SDEs is

\[
\boxed{
X_{n+1}
=
X_n
+
a(X_n,t_n)\Delta t
+
b(X_n,t_n)\Delta W_n,
}
\]

where

\[
\boxed{
\Delta W_n
\sim
N(0,\Delta t).
}
\]

This is the Euler--Maruyama scheme.

%%%%%%%%%%%%%%%%%%%%%%%%%%%%%%%%%%%%%%%%%%%%%%%%%%%%%%%%%%%%
\subsection{Monte Carlo Error and Time-Discretization Error}
\label{subsec:mc-sde-two-errors}
%%%%%%%%%%%%%%%%%%%%%%%%%%%%%%%%%%%%%%%%%%%%%%%%%%%%%%%%%%%%

SDE simulation introduces at least two numerical errors:

\[
\boxed{
\text{Monte Carlo sampling error}
}
\]

and

\[
\boxed{
\text{time-discretization error}.
}
\]

Thus increasing the number of paths cannot eliminate bias caused by a
coarse time discretization.

Both

\[
N
\]

and

\[
\Delta t
\]

must be controlled.

%%%%%%%%%%%%%%%%%%%%%%%%%%%%%%%%%%%%%%%%%%%%%%%%%%%%%%%%%%%%
\subsection{Strong and Weak SDE Accuracy Preview}
\label{subsec:strong-weak-sde}
%%%%%%%%%%%%%%%%%%%%%%%%%%%%%%%%%%%%%%%%%%%%%%%%%%%%%%%%%%%%

Strong convergence concerns approximation of individual sample paths, for
example

\[
\boxed{
\mathbb{E}
[
|X_T-X_T^{\Delta t}|
].
}
\]

Weak convergence concerns expectations:

\[
\boxed{
\left|
\mathbb{E}[g(X_T)]
-
\mathbb{E}[g(X_T^{\Delta t})]
\right|.
}
\]

Monte Carlo pricing and expectation estimation often care primarily about
weak accuracy.

%%%%%%%%%%%%%%%%%%%%%%%%%%%%%%%%%%%%%%%%%%%%%%%%%%%%%%%%%%%%
\subsection{Nested Monte Carlo}
\label{subsec:nested-monte-carlo}
%%%%%%%%%%%%%%%%%%%%%%%%%%%%%%%%%%%%%%%%%%%%%%%%%%%%%%%%%%%%

Some problems involve expectations inside expectations.

For example,

\[
\boxed{
\mathbb{E}
\left[
\phi
\left(
\mathbb{E}
[
Y
\mid
X
]
\right)
\right].
}
\]

An outer simulation samples \(X\).

For each outer sample, an inner simulation estimates the conditional
expectation.

This can be computationally expensive because total work is roughly the
product of outer and inner sample counts.

%%%%%%%%%%%%%%%%%%%%%%%%%%%%%%%%%%%%%%%%%%%%%%%%%%%%%%%%%%%%
\subsection{Uncertainty Quantification}
\label{subsec:monte-carlo-uq}
%%%%%%%%%%%%%%%%%%%%%%%%%%%%%%%%%%%%%%%%%%%%%%%%%%%%%%%%%%%%

Suppose a model output is

\[
Y
=
M(\Theta),
\]

where \(\Theta\) contains uncertain parameters.

Monte Carlo simulation propagates uncertainty through the model:

\[
\boxed{
\Theta_i
\longrightarrow
Y_i
=
M(\Theta_i).
}
\]

The output samples can estimate:

\[
\boxed{
\text{mean},
}
\]

\[
\boxed{
\text{variance},
}
\]

\[
\boxed{
\text{quantiles},
}
\]

and

\[
\boxed{
\text{tail probabilities}.
}
\]

%%%%%%%%%%%%%%%%%%%%%%%%%%%%%%%%%%%%%%%%%%%%%%%%%%%%%%%%%%%%
\subsection{Sensitivity Analysis With Monte Carlo}
\label{subsec:mc-sensitivity-analysis}
%%%%%%%%%%%%%%%%%%%%%%%%%%%%%%%%%%%%%%%%%%%%%%%%%%%%%%%%%%%%

Monte Carlo samples can be used to study how uncertainty in inputs affects
uncertainty in outputs.

Methods may examine:

\[
\boxed{
\text{correlations},
}
\]

\[
\boxed{
\text{variance decomposition},
}
\]

or

\[
\boxed{
\text{conditional output changes}.
}
\]

More advanced global sensitivity methods estimate the contribution of
individual parameters and their interactions.

%%%%%%%%%%%%%%%%%%%%%%%%%%%%%%%%%%%%%%%%%%%%%%%%%%%%%%%%%%%%
\subsection{Parallelism}
\label{subsec:mc-parallelism}
%%%%%%%%%%%%%%%%%%%%%%%%%%%%%%%%%%%%%%%%%%%%%%%%%%%%%%%%%%%%

Independent Monte Carlo samples can usually be computed independently.

Therefore simulation is naturally parallel:

\[
\boxed{
\text{sample }1
\parallel
\text{sample }2
\parallel
\cdots
\parallel
\text{sample }N.
}
\]

This makes Monte Carlo attractive on multicore, cluster, and distributed
systems.

%%%%%%%%%%%%%%%%%%%%%%%%%%%%%%%%%%%%%%%%%%%%%%%%%%%%%%%%%%%%
\subsection{Parallel Random Streams}
\label{subsec:parallel-random-streams}
%%%%%%%%%%%%%%%%%%%%%%%%%%%%%%%%%%%%%%%%%%%%%%%%%%%%%%%%%%%%

Parallel simulation must avoid unintended overlap between pseudorandom
streams.

Good implementations use properly designed independent or separated
streams.

Simply giving neighboring processors poorly chosen seeds may not guarantee
statistical independence.

%%%%%%%%%%%%%%%%%%%%%%%%%%%%%%%%%%%%%%%%%%%%%%%%%%%%%%%%%%%%
\subsection{Monte Carlo Convergence Study}
\label{subsec:mc-convergence-study}
%%%%%%%%%%%%%%%%%%%%%%%%%%%%%%%%%%%%%%%%%%%%%%%%%%%%%%%%%%%%

To verify the expected statistical rate, run simulations at increasing
sample sizes such as

\[
N,
\quad
4N,
\quad
16N.
\]

Because standard error scales approximately as

\[
N^{-1/2},
\]

multiplying \(N\) by \(4\) should reduce sampling error by about a factor of
\(2\).

%%%%%%%%%%%%%%%%%%%%%%%%%%%%%%%%%%%%%%%%%%%%%%%%%%%%%%%%%%%%
\subsection{Replicated Simulations}
\label{subsec:replicated-monte-carlo}
%%%%%%%%%%%%%%%%%%%%%%%%%%%%%%%%%%%%%%%%%%%%%%%%%%%%%%%%%%%%

Rather than performing one very large run, one may perform several
independent replications.

If estimates are

\[
\widehat{\mu}_1,
\ldots,
\widehat{\mu}_R,
\]

their variability provides empirical information about simulation error.

Replications are especially useful for checking theoretical standard-error
calculations.

%%%%%%%%%%%%%%%%%%%%%%%%%%%%%%%%%%%%%%%%%%%%%%%%%%%%%%%%%%%%
\subsection{Stopping by Desired Standard Error}
\label{subsec:mc-stopping-standard-error}
%%%%%%%%%%%%%%%%%%%%%%%%%%%%%%%%%%%%%%%%%%%%%%%%%%%%%%%%%%%%

Suppose the target standard error is

\[
\varepsilon.
\]

Since

\[
\operatorname{SE}
\approx
\frac{\sigma}{\sqrt{N}},
\]

a rough required sample size is

\[
\boxed{
N
\approx
\left(
\frac{
\sigma
}{
\varepsilon
}
\right)^2.
}
\]

A pilot simulation can estimate \(\sigma\).

%%%%%%%%%%%%%%%%%%%%%%%%%%%%%%%%%%%%%%%%%%%%%%%%%%%%%%%%%%%%
\subsection{Relative Monte Carlo Error}
\label{subsec:relative-mc-error}
%%%%%%%%%%%%%%%%%%%%%%%%%%%%%%%%%%%%%%%%%%%%%%%%%%%%%%%%%%%%

When the target mean \(\mu\neq0\), a relative Monte Carlo standard error is

\[
\boxed{
\frac{
\widehat{\operatorname{SE}}
}{
|\widehat{\mu}|
}.
}
\]

This is useful when quantities vary greatly in scale.

It becomes unstable when the true mean is close to zero.

%%%%%%%%%%%%%%%%%%%%%%%%%%%%%%%%%%%%%%%%%%%%%%%%%%%%%%%%%%%%
\subsection{Monte Carlo Workflow}
\label{subsec:monte-carlo-workflow}
%%%%%%%%%%%%%%%%%%%%%%%%%%%%%%%%%%%%%%%%%%%%%%%%%%%%%%%%%%%%

A practical workflow is:

\begin{enumerate}
    \item identify the quantity of interest;
    \item express it as an expectation when possible;
    \item determine the probability distribution to sample;
    \item choose an appropriate random-number generator;
    \item set and record seeds for reproducibility;
    \item generate a pilot sample;
    \item estimate variance;
    \item determine whether naive Monte Carlo is efficient;
    \item select variance-reduction methods if useful;
    \item determine sample size;
    \item generate samples;
    \item compute the estimator;
    \item estimate Monte Carlo standard error;
    \item construct uncertainty intervals;
    \item repeat using independent replications when useful;
    \item check convergence as \(N\) increases;
    \item inspect rare-event behavior;
    \item distinguish sampling error from discretization or model error;
    \item document the simulation configuration;
    \item verify results against known cases or alternative methods.
\end{enumerate}

%%%%%%%%%%%%%%%%%%%%%%%%%%%%%%%%%%%%%%%%%%%%%%%%%%%%%%%%%%%%
\subsection{Choosing a Monte Carlo Method}
\label{subsec:choosing-monte-carlo-method}
%%%%%%%%%%%%%%%%%%%%%%%%%%%%%%%%%%%%%%%%%%%%%%%%%%%%%%%%%%%%

For a directly simulatable expectation, use:

\[
\boxed{
\text{ordinary independent Monte Carlo}.
}
\]

For rare events, consider:

\[
\boxed{
\text{importance sampling}.
}
\]

For a correlated auxiliary variable with known mean, consider:

\[
\boxed{
\text{control variates}.
}
\]

For smooth integrals on a unit cube, consider:

\[
\boxed{
\text{quasi-Monte Carlo}.
}
\]

For a complicated target density available only up to normalization,
consider:

\[
\boxed{
\text{MCMC}.
}
\]

For time-evolving distributions, consider:

\[
\boxed{
\text{sequential Monte Carlo}.
}
\]

%%%%%%%%%%%%%%%%%%%%%%%%%%%%%%%%%%%%%%%%%%%%%%%%%%%%%%%%%%%%
\subsection{Common Errors}
\label{subsec:monte-carlo-common-errors}
%%%%%%%%%%%%%%%%%%%%%%%%%%%%%%%%%%%%%%%%%%%%%%%%%%%%%%%%%%%%

\subsubsection{Reporting an Estimate Without Its Monte Carlo Error}

A simulation output such as

\[
\widehat{\mu}=4.217
\]

is incomplete without information about uncertainty.

Report a standard error or confidence interval when possible.

%%%%%%%%%%%%%%%%%%%%%%%%%%%%%%%%%%%%%%%%%%%%%%%%%%%%%%%%%%%%
\subsubsection{Assuming More Samples Eliminate All Error}

Increasing \(N\) reduces sampling error.

It does not automatically reduce:

\[
\boxed{
\text{model bias},
}
\]

\[
\boxed{
\text{time-discretization bias},
}
\]

or

\[
\boxed{
\text{implementation error}.
}
\]

%%%%%%%%%%%%%%%%%%%%%%%%%%%%%%%%%%%%%%%%%%%%%%%%%%%%%%%%%%%%
\subsubsection{Confusing Standard Deviation and Standard Error}

The sample standard deviation estimates variation among individual
simulated outputs.

The standard error estimates variation of the sample mean:

\[
\boxed{
\operatorname{SE}
=
\frac{s}{\sqrt{N}}.
}
\]

%%%%%%%%%%%%%%%%%%%%%%%%%%%%%%%%%%%%%%%%%%%%%%%%%%%%%%%%%%%%
\subsubsection{Assuming Monte Carlo Converges Rapidly}

The standard rate

\[
O(N^{-1/2})
\]

is slow.

Doubling the number of samples does not halve the standard error.

%%%%%%%%%%%%%%%%%%%%%%%%%%%%%%%%%%%%%%%%%%%%%%%%%%%%%%%%%%%%
\subsubsection{Ignoring Random Seeds}

Without recording simulation settings and seeds, exact computational
reproduction may be difficult.

%%%%%%%%%%%%%%%%%%%%%%%%%%%%%%%%%%%%%%%%%%%%%%%%%%%%%%%%%%%%
\subsubsection{Treating Pseudorandom Numbers as Mathematically Independent Automatically}

Pseudorandom sequences are deterministic.

A high-quality generator is designed to approximate the statistical
behavior needed for simulation.

Generator quality matters.

%%%%%%%%%%%%%%%%%%%%%%%%%%%%%%%%%%%%%%%%%%%%%%%%%%%%%%%%%%%%
\subsubsection{Using a Poor Importance Distribution}

If

\[
q(x)
\]

is too small where

\[
g(x)p(x)
\]

is important, importance weights can become enormous.

Estimator variance may then explode.

%%%%%%%%%%%%%%%%%%%%%%%%%%%%%%%%%%%%%%%%%%%%%%%%%%%%%%%%%%%%
\subsubsection{Using Importance Sampling Without Support Coverage}

If

\[
q(x)=0
\]

where

\[
g(x)p(x)\neq0,
\]

the estimator cannot correctly represent that region.

The proposal must cover all relevant target support.

%%%%%%%%%%%%%%%%%%%%%%%%%%%%%%%%%%%%%%%%%%%%%%%%%%%%%%%%%%%%
\subsubsection{Assuming Antithetic Pairing Always Reduces Variance}

Antithetic variates help when the paired outputs are negatively
correlated.

Positive correlation may provide little benefit or even increase
variance.

%%%%%%%%%%%%%%%%%%%%%%%%%%%%%%%%%%%%%%%%%%%%%%%%%%%%%%%%%%%%
\subsubsection{Using a Weakly Correlated Control Variate}

Control variates are most effective when correlation magnitude is large.

A nearly unrelated control contributes little variance reduction.

%%%%%%%%%%%%%%%%%%%%%%%%%%%%%%%%%%%%%%%%%%%%%%%%%%%%%%%%%%%%
\subsubsection{Treating MCMC Samples as Independent}

Markov-chain samples are correlated.

Using

\[
s/\sqrt{N}
\]

as though all samples were independent can severely underestimate
uncertainty.

%%%%%%%%%%%%%%%%%%%%%%%%%%%%%%%%%%%%%%%%%%%%%%%%%%%%%%%%%%%%
\subsubsection{Assuming Burn-In Guarantees Convergence}

Discarding the first fixed number of samples does not prove that the chain
has reached the target distribution.

Mixing and convergence diagnostics are still needed.

%%%%%%%%%%%%%%%%%%%%%%%%%%%%%%%%%%%%%%%%%%%%%%%%%%%%%%%%%%%%
\subsubsection{Thinning Without a Reason}

Discarding most MCMC samples reduces stored autocorrelation but also throws
away information.

Thinning is not automatically required for valid MCMC estimation.

%%%%%%%%%%%%%%%%%%%%%%%%%%%%%%%%%%%%%%%%%%%%%%%%%%%%%%%%%%%%
\subsubsection{Judging MCMC Only by Acceptance Rate}

High acceptance can mean proposals are too small.

Low acceptance can mean proposals are too aggressive.

The important issue is efficient exploration of the target.

%%%%%%%%%%%%%%%%%%%%%%%%%%%%%%%%%%%%%%%%%%%%%%%%%%%%%%%%%%%%
\subsubsection{Using Too Few Chains for a Difficult Multimodal Target}

One chain can remain trapped in a single mode.

Multiple starting points may reveal unexplored regions.

%%%%%%%%%%%%%%%%%%%%%%%%%%%%%%%%%%%%%%%%%%%%%%%%%%%%%%%%%%%%
\subsubsection{Ignoring Rare Events}

A simulation may produce zero rare events simply because the sample size is
too small.

Zero observed events does not imply zero true probability.

%%%%%%%%%%%%%%%%%%%%%%%%%%%%%%%%%%%%%%%%%%%%%%%%%%%%%%%%%%%%
\subsubsection{Confusing Monte Carlo With Quasi-Monte Carlo}

Monte Carlo uses randomized samples.

Classical quasi-Monte Carlo uses deterministic low-discrepancy points.

Their error analyses are different.

%%%%%%%%%%%%%%%%%%%%%%%%%%%%%%%%%%%%%%%%%%%%%%%%%%%%%%%%%%%%
\subsubsection{Ignoring SDE Discretization Error}

Using millions of paths with a poor time discretization can produce a
highly precise estimate of a biased quantity.

Sampling and discretization errors must both be controlled.

%%%%%%%%%%%%%%%%%%%%%%%%%%%%%%%%%%%%%%%%%%%%%%%%%%%%%%%%%%%%
\subsection{Applications}
\label{subsec:monte-carlo-applications}
%%%%%%%%%%%%%%%%%%%%%%%%%%%%%%%%%%%%%%%%%%%%%%%%%%%%%%%%%%%%

\begin{applicationbox}

\textbf{Probability.}
Monte Carlo estimates probabilities, expectations, stopping-event
frequencies, and distributional quantities.

\medskip

\textbf{Statistics.}
Bootstrap methods, Bayesian posterior simulation, and randomization
procedures rely on repeated simulation.

\medskip

\textbf{Finance.}
Derivative pricing, portfolio risk, value-at-risk estimation, and
stochastic models frequently use Monte Carlo path simulation.

\medskip

\textbf{Physics.}
Statistical mechanics, particle systems, transport problems, and high-
dimensional integrals are major Monte Carlo applications.

\medskip

\textbf{Engineering.}
Reliability, uncertainty propagation, failure probability, and robust
design can be studied by random simulation.

\medskip

\textbf{Operations Research.}
Simulation evaluates queues, inventories, service systems, and stochastic
decision policies.

\medskip

\textbf{Bayesian Inference.}
MCMC and related methods approximate posterior expectations when analytic
integration is impossible.

\medskip

\textbf{Scientific Computing.}
Monte Carlo provides a general computational framework for problems where
random sampling is easier than deterministic integration.

\end{applicationbox}

%%%%%%%%%%%%%%%%%%%%%%%%%%%%%%%%%%%%%%%%%%%%%%%%%%%%%%%%%%%%
\subsection{Exercises}
\label{subsec:monte-carlo-exercises}
%%%%%%%%%%%%%%%%%%%%%%%%%%%%%%%%%%%%%%%%%%%%%%%%%%%%%%%%%%%%

\begin{exercise}[1]
Define a Monte Carlo estimator.
\end{exercise}

\begin{exercise}[1]
State the standard Monte Carlo convergence rate.
\end{exercise}

\begin{exercise}[1]
Define Monte Carlo standard error.
\end{exercise}

\begin{exercise}[1]
Define importance sampling.
\end{exercise}

\begin{exercise}[1]
Define a control variate.
\end{exercise}

\begin{exercise}[1]
Define MCMC.
\end{exercise}

\begin{exercise}[1]
Define a stationary distribution.
\end{exercise}

\begin{exercise}[1]
Define effective sample size conceptually.
\end{exercise}

\begin{exercise}[1]
Define quasi-Monte Carlo conceptually.
\end{exercise}

\begin{exercise}[1]
Define a bootstrap sample.
\end{exercise}

\begin{exercise}[2]
Compute the Monte Carlo estimate from a supplied sample.
\end{exercise}

\begin{exercise}[2]
Compute the sample variance of simulated outputs.
\end{exercise}

\begin{exercise}[2]
Compute the Monte Carlo standard error.
\end{exercise}

\begin{exercise}[2]
Construct an approximate \(95\%\) confidence interval.
\end{exercise}

\begin{exercise}[2]
Determine how the standard error changes if \(N\) is multiplied by \(4\).
\end{exercise}

\begin{exercise}[2]
Determine how much \(N\) must increase to reduce standard error by a
factor of \(5\).
\end{exercise}

\begin{exercise}[2]
Write

\[
\int_0^2f(x)\,dx
\]

as an expectation involving a uniform random variable.
\end{exercise}

\begin{exercise}[3]
Construct a Monte Carlo estimator for that integral.
\end{exercise}

\begin{exercise}[3]
Estimate \(\pi\) from a supplied set of random points.
\end{exercise}

\begin{exercise}[3]
Derive the variance of a Monte Carlo probability estimator.
\end{exercise}

\begin{exercise}[3]
Estimate the number of samples needed to achieve a desired standard error.
\end{exercise}

\begin{exercise}[2]
Explain the purpose of a random seed.
\end{exercise}

\begin{exercise}[2]
Use inverse-transform sampling to simulate an exponential random variable.
\end{exercise}

\begin{exercise}[3]
Derive the inverse-transform formula for a specified continuous
distribution.
\end{exercise}

\begin{exercise}[3]
Design an acceptance--rejection sampler for a target density with a given
proposal.
\end{exercise}

\begin{exercise}[3]
Compute the acceptance probability when the bounding constant \(M\) is
given.
\end{exercise}

\begin{exercise}[2]
Derive the basic importance-sampling identity.
\end{exercise}

\begin{exercise}[3]
Construct an importance-sampling estimator for a specified integral.
\end{exercise}

\begin{exercise}[3]
Explain why support coverage is necessary.
\end{exercise}

\begin{exercise}[3]
Compute importance weights for given target and proposal densities.
\end{exercise}

\begin{exercise}[3]
Explain why importance sampling can help estimate rare-event
probabilities.
\end{exercise}

\begin{exercise}[3]
Derive the relative standard error of a naive rare-event probability
estimator.
\end{exercise}

\begin{exercise}[2]
Construct a control-variate estimator.
\end{exercise}

\begin{exercise}[3]
Derive the optimal control-variate coefficient.
\end{exercise}

\begin{exercise}[3]
Show that the resulting variance depends on

\[
1-\rho^2.
\]
\end{exercise}

\begin{exercise}[2]
Construct an antithetic estimator using \(U\) and \(1-U\).
\end{exercise}

\begin{exercise}[3]
Explain when antithetic pairing reduces variance.
\end{exercise}

\begin{exercise}[3]
Construct a two-stratum Monte Carlo estimator.
\end{exercise}

\begin{exercise}[3]
Explain Latin hypercube sampling.
\end{exercise}

\begin{exercise}[3]
Explain common random numbers for comparing two systems.
\end{exercise}

\begin{exercise}[2]
Write the sample-average approximation of

\[
F(x)
=
\mathbb{E}[f(x,\xi)].
\]
\end{exercise}

\begin{exercise}[3]
Construct a Monte Carlo gradient estimator.
\end{exercise}

\begin{exercise}[2]
Write the Metropolis--Hastings acceptance probability.
\end{exercise}

\begin{exercise}[3]
Compute a Metropolis--Hastings acceptance probability for supplied
densities.
\end{exercise}

\begin{exercise}[3]
Simplify the acceptance rule for a symmetric proposal.
\end{exercise}

\begin{exercise}[3]
Explain why an unknown normalizing constant cancels.
\end{exercise}

\begin{exercise}[3]
Explain detailed balance.
\end{exercise}

\begin{exercise}[3]
Describe one Gibbs sampling sweep for a bivariate distribution.
\end{exercise}

\begin{exercise}[3]
Explain why MCMC samples are not independent.
\end{exercise}

\begin{exercise}[3]
Compute effective sample size from a supplied integrated autocorrelation
time.
\end{exercise}

\begin{exercise}[3]
Explain the purpose and limitation of burn-in.
\end{exercise}

\begin{exercise}[3]
Explain why thinning is not automatically necessary.
\end{exercise}

\begin{exercise}[3]
Describe indicators of poor MCMC mixing.
\end{exercise}

\begin{exercise}[3]
Explain why multiple chains can be useful.
\end{exercise}

\begin{exercise}[3]
Construct a nonparametric bootstrap sample from five observations.
\end{exercise}

\begin{exercise}[3]
Estimate a bootstrap standard error from supplied bootstrap statistics.
\end{exercise}

\begin{exercise}[3]
Distinguish parametric and nonparametric bootstrap.
\end{exercise}

\begin{exercise}[3]
Explain low-discrepancy sampling.
\end{exercise}

\begin{exercise}[3]
Compare ordinary Monte Carlo and quasi-Monte Carlo.
\end{exercise}

\begin{exercise}[3]
Write the Euler--Maruyama scheme.
\end{exercise}

\begin{exercise}[3]
Simulate one Euler--Maruyama time step for supplied values of
\(\Delta W\).
\end{exercise}

\begin{exercise}[3]
Explain the difference between strong and weak SDE accuracy.
\end{exercise}

\begin{exercise}[3]
Explain why increasing the number of SDE paths does not remove time-step
bias.
\end{exercise}

\begin{exercise}[3]
Design a Monte Carlo convergence study.
\end{exercise}

\begin{exercise}[3]
Estimate the empirical rate at which simulation standard error decreases.
\end{exercise}

\begin{exercise}[3]
Explain why Monte Carlo computations are naturally parallel.
\end{exercise}

\begin{exercise}[4]
Prove unbiasedness of the sample-average Monte Carlo estimator.
\end{exercise}

\begin{exercise}[4]
Derive the variance

\[
\operatorname{Var}
(
\widehat{\mu}_N
)
=
\frac{\sigma^2}{N}.
\]
\end{exercise}

\begin{exercise}[4]
Derive the central-limit confidence interval for a Monte Carlo mean.
\end{exercise}

\begin{exercise}[4]
Study Berry--Esseen corrections to Monte Carlo normal approximation.
\end{exercise}

\begin{exercise}[4]
Study optimal allocation in stratified sampling.
\end{exercise}

\begin{exercise}[4]
Derive the minimum-variance importance-sampling density formally.
\end{exercise}

\begin{exercise}[4]
Study self-normalized importance sampling.
\end{exercise}

\begin{exercise}[4]
Study exponential tilting for rare events.
\end{exercise}

\begin{exercise}[4]
Study conditional Monte Carlo variance reduction.
\end{exercise}

\begin{exercise}[4]
Study Rao--Blackwellization as variance reduction.
\end{exercise}

\begin{exercise}[4]
Prove the optimal control-variate coefficient.
\end{exercise}

\begin{exercise}[4]
Study Latin hypercube variance properties.
\end{exercise}

\begin{exercise}[4]
Study randomized quasi-Monte Carlo.
\end{exercise}

\begin{exercise}[4]
Study discrepancy and the Koksma--Hlawka inequality.
\end{exercise}

\begin{exercise}[4]
Prove detailed balance for Metropolis--Hastings.
\end{exercise}

\begin{exercise}[4]
Study irreducibility and aperiodicity in MCMC.
\end{exercise}

\begin{exercise}[4]
Study Markov-chain ergodic theorems.
\end{exercise}

\begin{exercise}[4]
Derive MCMC asymptotic variance.
\end{exercise}

\begin{exercise}[4]
Study effective sample size estimators.
\end{exercise}

\begin{exercise}[4]
Study adaptive MCMC conceptually.
\end{exercise}

\begin{exercise}[4]
Study Hamiltonian Monte Carlo conceptually.
\end{exercise}

\begin{exercise}[4]
Study particle filters and sequential Monte Carlo.
\end{exercise}

\begin{exercise}[4]
Study bootstrap consistency.
\end{exercise}

\begin{exercise}[4]
Study bias-corrected bootstrap intervals.
\end{exercise}

\begin{exercise}[4]
Analyze weak convergence of Euler--Maruyama.
\end{exercise}

\begin{exercise}[4]
Study multilevel Monte Carlo.
\end{exercise}

\begin{exercise}[4]
Investigate nested Monte Carlo complexity.
\end{exercise}

%%%%%%%%%%%%%%%%%%%%%%%%%%%%%%%%%%%%%%%%%%%%%%%%%%%%%%%%%%%%
\subsection{Conceptual Summary}
\label{subsec:monte-carlo-summary}
%%%%%%%%%%%%%%%%%%%%%%%%%%%%%%%%%%%%%%%%%%%%%%%%%%%%%%%%%%%%

Monte Carlo methods approximate expectations using random samples.

For

\[
\boxed{
\mu
=
\mathbb{E}[g(X)],
}
\]

the standard estimator is

\[
\boxed{
\widehat{\mu}_N
=
\frac1N
\sum_{i=1}^{N}
g(X_i).
}
\]

Its variance is

\[
\boxed{
\operatorname{Var}
(
\widehat{\mu}_N
)
=
\frac{
\sigma^2
}{
N
},
}
\]

so its standard error is

\[
\boxed{
O(N^{-1/2}).
}
\]

Monte Carlo integration rewrites an integral as an expectation.

Variance-reduction methods include

\[
\boxed{
\text{importance sampling}
+
\text{control variates}
+
\text{antithetic variates}
+
\text{stratification}.
}
\]

When direct independent sampling is difficult, MCMC constructs a Markov
chain whose stationary distribution is the target.

Metropolis--Hastings uses acceptance probability

\[
\boxed{
\alpha(x,y)
=
\min
\left\{
1,
\frac{
\pi(y)q(x\mid y)
}{
\pi(x)q(y\mid x)
}
\right\}.
}
\]

Monte Carlo methods therefore combine

\[
\boxed{
\text{probability}
+
\text{statistics}
+
\text{numerical approximation}
+
\text{random simulation}
+
\text{computational inference}.
}
\]

%%%%%%%%%%%%%%%%%%%%%%%%%%%%%%%%%%%%%%%%%%%%%%%%%%%%%%%%%%%%
\subsection{Section Connections}
\label{subsec:monte-carlo-connections}
%%%%%%%%%%%%%%%%%%%%%%%%%%%%%%%%%%%%%%%%%%%%%%%%%%%%%%%%%%%%

\chapterconnection{
Monte Carlo Methods connect the probabilistic foundations of Chapter~7
with the statistical methods of Chapter~8 and the computational framework
of Chapter~11. Ordinary Monte Carlo converts expectations and integrals
into sample averages, while importance sampling and control variates
improve efficiency. MCMC provides a computational route to complicated
probability distributions, especially in Bayesian statistics, and
Euler--Maruyama connects simulation with the stochastic differential
equations of Chapter~7. Monte Carlo also appears in stochastic
optimization, reliability analysis, finance, uncertainty quantification,
and large-scale scientific simulation. The next section, Scientific
Computing, broadens the computational viewpoint to algorithm
implementation, software architecture, precision, performance,
parallelism, reproducibility, and large-scale mathematical computation.
}

%%%%%%%%%%%%%%%%%%%%%%%%%%%%%%%%%%%%%%%%%%%%%%%%%%%%%%%%%%%%
% SECTION SOURCE TRACKING
%%%%%%%%%%%%%%%%%%%%%%%%%%%%%%%%%%%%%%%%%%%%%%%%%%%%%%%%%%%%

%------------------------------------------------------------
% 11.6 Monte Carlo Methods
%------------------------------------------------------------
%
% Source basis:
%   - Standard Monte Carlo simulation theory.
%   - Standard variance-reduction methods.
%   - Standard Markov chain Monte Carlo.
%   - Standard quasi-Monte Carlo, bootstrap, and stochastic simulation
%     concepts.
%
% Used for:
%   - expectation estimation;
%   - Monte Carlo sample averages;
%   - unbiasedness;
%   - estimator variance;
%   - N^{-1/2} convergence;
%   - Law of Large Numbers;
%   - Central Limit Theorem;
%   - Monte Carlo standard errors;
%   - confidence intervals;
%   - bias and variance;
%   - Monte Carlo integration;
%   - high-dimensional integration;
%   - area and volume estimation;
%   - probability estimation;
%   - indicator variables;
%   - pseudorandom number generation;
%   - random seeds;
%   - inverse-transform sampling;
%   - transformation methods;
%   - acceptance--rejection sampling;
%   - variance reduction;
%   - importance sampling;
%   - importance weights;
%   - rare-event simulation;
%   - control variates;
%   - antithetic variates;
%   - stratified sampling;
%   - Latin hypercube sampling preview;
%   - common random numbers;
%   - Monte Carlo optimization;
%   - sample-average approximation;
%   - gradient estimation;
%   - Markov Chain Monte Carlo;
%   - stationary distributions;
%   - Metropolis--Hastings;
%   - symmetric proposals;
%   - unnormalized targets;
%   - detailed balance;
%   - Gibbs sampling;
%   - burn-in;
%   - autocorrelation;
%   - effective sample size;
%   - mixing and diagnostics;
%   - multiple chains;
%   - sequential Monte Carlo preview;
%   - particle degeneracy;
%   - bootstrap simulation;
%   - parametric and nonparametric bootstrap;
%   - quasi-Monte Carlo;
%   - discrepancy;
%   - randomized quasi-Monte Carlo;
%   - stochastic differential equation simulation;
%   - Euler--Maruyama;
%   - strong and weak SDE accuracy preview;
%   - nested Monte Carlo;
%   - uncertainty quantification;
%   - sensitivity analysis;
%   - parallel simulation;
%   - convergence studies;
%   - applications and exercises.
%
% Notes:
%   - Scientific Computing is developed next in Section 11.7.
%   - Stochastic-process foundations appear in Chapter 7.
%   - Statistical inference appears in Chapter 8.
%   - Stochastic Optimization appears in Section 10.8.
%
%%%%%%%%%%%%%%%%%%%%%%%%%%%%%%%%%%%%%%%%%%%%%%%%%%%%%%%%%%%%